% Generated from validated Article Draft Result. Do not hand-edit; revise the structured draft instead.
\section{Agents, Runtime \& Long-horizon Work — モデルを実行環境へ置く}
\noindent\textbf{Responses APIのcomputer environment、long-horizon reliability、deep research verification、agent memory。3月には、Agentを単なるtool callではなく、shell・container・file・stateと長時間評価を伴うsystemとして捉える材料が揃い始めた。}\autocite{src-01572b4c34e9,src-0ec537ec1cc1,src-2b8b39e5cae3,src-42318ebdee75}

\subsection{Agentは最大テーマではない、それでも構造変化は見える}

3月のCandidate surfaceではAgent/Runtime単独の件数はSystemsやSecurityを下回る。したがって今月全体を「Agentの月」と呼ぶのは強すぎる。一方で、Responses APIのcomputer environment、long-horizon reliability、verification-centric deep research、memory frameworkという材料は、modelを実行環境・状態・評価と結びつける方向が具体化していたことを示す。\autocite{src-01572b4c34e9,src-0ec537ec1cc1,src-2b8b39e5cae3,src-42318ebdee75}


\subsection{Responses API: shell・container・files・tools・stateを持つruntime}

OpenAIは3月11日、Responses API、shell tool、hosted containersを組み合わせ、files、tools、stateを持つagent runtimeを構築したと説明した。ここで変わるのはmodelの推論能力そのものではなく、推論結果を実際のcomputer environmentで継続的なactionへ変換する周辺systemである。\autocite{src-01572b4c34e9}


\subsection{Long-horizonではpass@1だけでは足りない}

Beyond pass@1はauthorsが、task durationが伸びるにつれてreliabilityの性質が分岐し、短時間taskのpass@1ではその差を捉えにくいと主張する。Agentを長く動かすなら、一回の成功率だけでなく、途中失敗、回復、累積的な安定性をどう評価するかが別の問題になる。\autocite{src-0ec537ec1cc1}


\subsection{Memoryとverificationをruntimeの一部として扱う}

Marco DeepResearchはauthorsがQA data synthesisやtrajectory constructionなど複数段階へverificationを組み込むdeep research designを提案した。MemFactoryはmemory-augmented agents向けにtrainingとinferenceを統合するmodular frameworkを提示している。長時間Agentでは「何を覚えるか」と「途中結果をどう確かめるか」がmodel外部の設計変数になる。\autocite{src-2b8b39e5cae3,src-42318ebdee75}


この4件だけで3月をAgent中心の転換点と断定する必要はない。ただし、execution environment、long-horizon evaluation、verification、memoryという構成要素が同じ月に並んだことで、model capabilityをsystemとして運用するための論点が輪郭を持ち始めていた、と位置付けることはできる。\autocite{src-01572b4c34e9,src-0ec537ec1cc1,src-2b8b39e5cae3,src-42318ebdee75}


\begin{claimboundary}
Claim Boundary: Responses API runtimeの安全性・scalabilityはOpenAIのfirst-party説明であり、本誌が独立検証したものではない。long-horizon、deep research、memory frameworkの結果はauthorsによる報告である。MemFactoryは現在v4 locatorを参照するため、March chronologyは初回投稿を基準にし、後日revisionの内容を3月へ遡及させない。\autocite{src-01572b4c34e9,src-0ec537ec1cc1,src-2b8b39e5cae3,src-42318ebdee75}
\end{claimboundary}
